
% ════════════════════════════════════════
% 结论（12.结论.tex）
% ════════════════════════════════════════
\section{结论}

本文围绕微网与外部电网之间的电力调控问题，建立了从确定性线性规划到两阶段随机规划、再到滚动时域调整的统一优化框架，并在固定电价与波动电价两种场景下逐一求解，主要结论如下。

问题一是确定性场景下的最优购电与储能调度。在功率平衡、储能容量与日循环约束下，模型以全天购电费最小为目标构成线性规划，求得全天计划购电量 59482.70 kWh、购电费 35126.95 元，储能于 0:00 与 24:00 的储电量均为附录1给定的 6000 kWh，并在低价时段充电、高价时段放电完成日内套利。对照检验表明，若取消储能，全天购电费将升至 48052.05 元、较配置储能时高出 36.80\%，可见削峰填谷是购电费用下降的主要来源。

问题二在预测不确定性下给出全年逐日滚动策略。模型先以扩张窗口标定负载与光伏预测、再由历史残差构造偏差情景，基于两阶段随机规划制定当日计划，最后按 5 倍电价结算紧急购电与弃电。2 月 1 日至 12 月 31 日共 334 天，计划购电费 13325494.59 元、紧急购电费 1057398.40 元，全年总购电费 14382892.99 元，其中紧急购电量 282910.68 kWh、弃电量 2881875.41 kWh。能量平衡校验显示，计划段与紧急段之间的收支恰好闭合，说明逐日结算无遗漏；弃电主要来自午间光伏超发，紧急购电则集中于夜间与傍晚光伏不足的时段。

问题三在计划可改、改价分段的前提下建立含违约惩罚的滚动调整模型。用上境变量将分段凸费用精确线性化后，模型仍是线性规划，可求得全局最优解。结果显示，在 0:00 计划的基础上再引入 6:00、12:00、18:00 三次预报制定调整策略，可节省购电费 444073.36 元、约 3.00\%，即全年总费用降至 14337999.15 元。进一步的边际分析表明，6 小时间隔已几乎捕获全部频率收益，继续加密至 3 小时或 1 小时仅再节省 5.12 万元与 0.05 万元，因此若需继续降低成本，应优先提高预报精度而非增加预报时刻。

问题四考察实时波动电价下的策略。将固定电价序列替换为附件四的逐间隔波动电价后重新求解，问题二的总费用升至 15077131.66 元，问题三降至 14956526.18 元、较问题二节省 120605.48 元（约 0.80\%）。这说明波动电价抬高了微网购电的不确定性与成本，而滚动调整与储能套利只能部分对冲电价波动的影响，模型的相对结论在两种电价机制下保持一致。

在方法层面，本文的贡献可归纳为三点：一是提出“先预测、再决策、后结算”的三段式框架，把预测、随机决策与事后结算解耦，使计划量不含事后信息；二是用上境变量将“计划可改、改价分段”的凸费用问题线性化，在不引入整数变量的前提下精确刻画违约与超额的经济代价；三是通过逐日滚动与储电量传递，将全年大规模优化分解为数百个日尺度小问题，在保证全局最优性的同时兼顾计算效率。

本文仍存在若干局限：预测误差只以三档情景近似，未显式建模其相关性，对极端天气的刻画偏粗；午间弃电未被进一步利用；储能效率取为单一常数，未计入循环寿命损耗。这些方面均可作为后续改进方向，相关推广已在第 8 节讨论。
